% BHU PYQ -- History: The Art of History Writing (2025-26 NEP)
\documentclass[11pt,a4paper]{article}
\usepackage[a4paper,top=2.5cm,bottom=2.5cm,left=2.8cm,right=2.8cm,headheight=14pt]{geometry}
\usepackage{amsmath,amssymb}
\usepackage{fontspec}
\setmainfont{Kohinoor Devanagari}
\usepackage{microtype,fancyhdr,enumitem,hyperref,lastpage,parskip}

\hypersetup{colorlinks=true,urlcolor=black,linkcolor=black,pdftitle={HISMV31: The Art of History Writing -- BHU NEP 2025-26}}

\pagestyle{fancy}\fancyhf{}
\renewcommand{\headrulewidth}{0.4pt}\renewcommand{\footrulewidth}{0.4pt}
\fancyhead[L]{\small \textit{Banaras Hindu University}}
\fancyhead[R]{\small \textit{HISMV31: Art of History Writing (2025-26)}}
\fancyfoot[C]{\small Page \thepage\ of \pageref{LastPage}}

\newlist{parts}{enumerate}{2}
\setlist[parts,1]{label=(\alph*),leftmargin=*,itemsep=4pt,topsep=3pt}
\setlist[parts,2]{label=(\roman*),leftmargin=*,itemsep=2pt,topsep=2pt}

\newcommand{\pts}[1]{\hfill \textit{[#1]}}

\begin{document}

\begin{center}
  {\large\bfseries BANARAS HINDU UNIVERSITY}\\[2pt]
  {\normalsize B. A. (Hons.) Semester III Examination, 2025-26 (NEP)}\\[6pt]
  \rule{0.8\linewidth}{0.5pt}\\[6pt]
  {\large\bfseries HISTORY}\\[4pt]
  {\large\bfseries Paper Code: HISMV-31 : The Art of History Writing}\\[4pt]
  {\normalsize Time: 2:30 Hours \hfill Full Marks: 70}\\[2pt]
  {\small REF NO. 2025-26/D-III/143}
\end{center}
\rule{\linewidth}{0.4pt}
\smallskip

\noindent \textit{Instructions: This question paper comprises three Sections A, B and C. All questions are compulsory. Write your Roll No.\ at the top immediately on receipt of this question paper.}

\smallskip
\noindent \textit{इस प्रश्न-पत्र में तीन खण्ड अ, ब तथा स हैं। सभी प्रश्न अनिवार्य हैं।}

\bigskip

\section*{SECTION -- A / खण्ड -- अ}
\noindent \textit{Note: Very Short Answer Type Questions. Answer each question.}\pts{5 $\times$ 2 = 10}

\noindent \textit{अति लघु उत्तरीय प्रश्न। प्रत्येक प्रश्न 2 अंकों का है।}

\begin{enumerate}[label=\textbf{\arabic*.},leftmargin=*,itemsep=8pt]
  \item 
  \begin{parts}
    \item What is a book?\pts{2}
    
    किताब (पुस्तक) क्या है?

    \item What is the difference between alphabets and script?\pts{2}
    
    वर्णमाला (अक्षर) तथा लिपि के मध्य क्या अन्तर है?

    \item What is an electronic citation?\pts{2}
    
    इलेक्ट्रॉनिक साइटेशन (ई-उद्धरण) क्या है?

    \item What do you know about annotated bibliography?\pts{2}
    
    समीक्षात्मक (व्याख्यात्मक) ग्रन्थ-सूची से आप क्या समझते हैं?

    \item Which was the first book to be printed in the modern age?\pts{2}
    
    आधुनिक युग में मुद्रित होने वाली पहली पुस्तक कौन-सी थी? (Gutenberg Bible)
  \end{parts}
\end{enumerate}

\bigskip

\section*{SECTION -- B / खण्ड -- ब}
\noindent \textit{Note: Short Answer Type Questions. Answer each question.}\pts{3 $\times$ 10 = 30}

\noindent \textit{लघु उत्तरीय प्रश्न। प्रत्येक प्रश्न 10 अंकों का है।}

\begin{enumerate}[label=\textbf{\arabic*.},leftmargin=*,start=2,itemsep=12pt]

  \item Define what is a Blurb. Throw light on its utility value.\pts{10}
  
  `ब्लर्ब' (Blurb) क्या है? इसकी उपयोगिता एवं मूल्य पर प्रकाश डालिए।

  \begin{center}\textbf{OR / अथवा}\end{center}

  Give a short account of the forms of book reviews and their scope.\pts{10}
  
  पुस्तक समीक्षाओं के रूपों तथा उनके क्षेत्र का संक्षिप्त विवरण दीजिए।

  \item Critically analyze the steps involved in the process of making a book.\pts{10}
  
  पुस्तक निर्माण की प्रक्रिया में शामिल चरणों का समालोचनात्मक विश्लेषण कीजिए।

  \begin{center}\textbf{OR / अथवा}\end{center}

  What is desktop publishing (DTP) and its market value?\pts{10}
  
  डेस्कटॉप पब्लिशिंग (DTP) क्या है तथा इसका बाजार मूल्य क्या है?

  \item What are historical sources? How have they been used in history writing? Explain with examples.\pts{10}
  
  ऐतिहासिक स्रोत क्या हैं? इतिहास लेखन में इनका उपयोग किस प्रकार किया गया है? उदाहरण सहित समझाइए।

  \begin{center}\textbf{OR / अथवा}\end{center}

  Trace the development of printing technology since ancient times.\pts{10}
  
  प्राचीन काल से मुद्रण प्रौद्योगिकी के विकास को रेखांकित कीजिए।

\end{enumerate}

\bigskip

\section*{SECTION -- C / खण्ड -- स}
\noindent \textit{Note: Long Answer Type Questions. Answer each question.}\pts{2 $\times$ 15 = 30}

\noindent \textit{दीर्घ उत्तरीय प्रश्न। प्रत्येक प्रश्न 15 अंकों का है।}

\begin{enumerate}[label=\textbf{\arabic*.},leftmargin=*,start=5,itemsep=14pt]

  \item Critically analyze the steps involved in reviewing a book by using the example of Barrington Moore's \textit{``Social Origins of Dictatorship and Democracy''}.\pts{15}
  
  बैरिंगटन मूरे की पुस्तक ``सोशल ओरिजिन्स ऑफ डिक्टेटरशिप एंड डेमोक्रेसी'' का उदाहरण लेते हुए पुस्तक समीक्षा के विभिन्न चरणों का समालोचनात्मक विश्लेषण कीजिए।

  \begin{center}\textbf{OR / अथवा}\end{center}

  Critically analyze the process of improving historical writing with various modern writing software applications.\pts{15}
  
  विभिन्न सॉफ्टवेयर एप्लिकेशनों के माध्यम से अपने ऐतिहासिक लेखन में सुधार की प्रक्रिया का समालोचनात्मक विश्लेषण कीजिए।

  \item ``The pen is mightier than the sword.'' How far is this statement true in the current historical context?\pts{15}
  
  ``कलम तलवार से अधिक शक्तिशाली होती है।'' वर्तमान ऐतिहासिक संदर्भ में यह कथन कहाँ तक सत्य है?

  \begin{center}\textbf{OR / अथवा}\end{center}

  Enumerate the different forms of literary genres in historical writing with appropriate examples.\pts{15}
  
  ऐतिहासिक लेखन में साहित्यिक विधाओं के विभिन्न रूपों का उदाहरण सहित विवरण दीजिए।

\end{enumerate}

\end{document}
